\section{Departures of the implementation and boundaries of the text}
\label{sec:departures}

The extended edition keeps a ledger of proof obligations open against the
current specification, CCME 3.1.0-rc.2 at commit \code{252216ba}, and against
the reference implementation at the same commit, each with the scenario that
exposes it, its consequence for the results of this report, the repair or the
recorded status, and the vector and test that pin it. \Cref{tab:departures}
states each in one row; the F-numbered references open the full entries. The
sources are \SPEC{} itself, its dated design history and changelog, the
engine's release notes, the integration test plan, and the test sources of
the engine and its integration suite. The six obligations found in CCME 1.0.0
(the prerelease window that never discharged, the suppression lapse, the
staleness dual without the resolvability test, the one-directional axis
independence, the inherited channel that read a moving baseline, and the
bucketing remedy's own edge) were repaired in CCME 2.0.0 and are reproduced
in the extended edition as provenance with their minimal counterexamples and
regression vectors (\cref{ob:pre,ob:g5,ob:duality,ob:g8,ob:origin,ob:bucket}).

\begin{center}
\small
\begin{longtable}{@{}L{2.2cm}L{6.2cm}L{5.6cm}@{}}
\caption{Current obligations, ordered by severity: the silent loss of a
release first, diagnostic metadata last.}
\label{tab:departures}\\
\toprule
Entry & Claim & Status and consequence \\
\midrule
\endfirsthead
\toprule
Entry & Claim & Status and consequence \\
\midrule
\endhead
\Cref{ob:owed} & \spec{13.3} adds an owed window per dependency edge and
\spec{19.3} refuses with \diag{E201} a run that would release a provider at
the baseline commit of a consumer it still owes. The engine implements the
delivery test and the reconciliation of a proceeding consumer, and neither the
owed window nor \diag{E201}. & Open; the engine's own design history records
it. A consumer that sits out the run in which its provider delivers is
stranded: \cref{thm:g2} holds in the specification and, in the engine, only
while some pending window still holds the commit. The loss is silent in the
plan. Pinned by vector 80d; no test exercises the run the consumer sits out. \\
\Cref{ob:attribution} & In the builds the campaigns ran, the bump axis credited
a unit's whole source set to every dependent the walk reached, where
\spec{9.2} attributes only the sources within the unit's depth. & Fixed on the
product repository's \code{main} at \code{e461d7d9}, after the measured builds
and unpublished when this was written. No effect on \cref{sec:eval:rq5}: every
trigger commit there scopes one package, for which both rules attribute
identically. It is the (R) shape of \cref{sec:guarantees}. \\
\Cref{ob:samecommit} & \spec{19.3}: if the consumer is published after the
provider and fails, two releases sit on one commit and no later plan can
compute the debt; \diag{E201} reports it with its remedy, an explicit
\code{Release-As}. & A boundary the specification states. Completeness
(\cref{thm:g2}) reports this state rather than resolving it; the engine, which
lacks \diag{E201}, does not report it. \\
\Cref{ob:graduation} & \spec{13.9}'s stable branch computes the graduation
version from \code{effective(P)} over surviving tuples, while \spec{11.5}
defines it as the train target accumulated over the pending window. & A
contradiction inside the text, unchanged since 2.0.0. \Cref{thm:g3} is proved
for fixed version inputs; which of the two rules an implementation follows is
pinned by the last row of \spec{11.4}'s table. \\
\Cref{ob:kinds} & The engine builds one dependency set and uses it for
propagation, ordering, blocking and the cycle check; the text distinguishes
\code{propagation.kinds}, \code{publish.orderKinds} and
\code{publish.blockingKinds}, and excludes a dev-only cycle from \diag{E200}. &
Open, on the safe side: the engine over-blocks and refuses dev-only cycles it
should accept. The theorems assume hypothesis (K), which the engine satisfies
trivially. Pinned by vectors 79 and 100a and edge cases 99a and 100. \\
\Cref{ob:records} & \spec{19.1}: a release tag is created, never replaced,
except for the retry of an uncertain write that names the same commit,
package and version. & The single-repository recorder implements both halves;
the fleet recorder of \spec{27} implements the rule and not the exception, so
a retry after an applied-but-unacknowledged push is refused rather than
recognised. Nothing incorrect is recorded. Open for the fleet path. \\
\Cref{ob:category} & \spec{28.9} requires a stable machine-readable failure
class on every listed row. & The class is attached where a distributed run
loses its lock and not yet at the older recording and lock sites; the engine's
design history records it as a departure to be closed. Pinned by vector 14 of
\spec{28.8}. \\
\Cref{ob:profile} & The design history lists what the engine does not
implement of \spec{28} and where it departs: reuse of verified outputs across
runs (vector 20) and the rollback profile of \spec{26} (vector 26) are not
implemented; the platform check runs before dispatch for released packages
and at placement for a prepared provider; the summary shape of a prepared
provider and the phase a cancelled attempt stopped in were added to the text
from the implementation. & Recorded so that the conformance table of the
extended edition reads against them; none affects a proved theorem. \\
\bottomrule
\end{longtable}
\end{center}

Beside these, the extended edition records textual defects noticed in passing
that change no result, and the dates and commits of the design-history entry
of 22 September 2026 as observed while the product side corrected them.
